\bookchapter{제29장: 디지털 건강——스마트폰과 건강한 관계 구축}

\chapterintro{5분만 숏폼 영상 보려고 했는데, 고개를 들어보니 이미 두 시간이 지나 있었다는 경험이 있나요? 또는 밤에 누워서 스마트폰을 보면서, 볼수록 더 정신이 깨어나, 결국 새벽 2시에야 잠드는 경험. 스마트폰은 우리의 삶을 바꾸었지만, 또한 조용히 우리의 시간과 건강을 잠식하고 있습니다. 이 챕터는 “스마트폰을 끊으라”가 아니라, 스마트폰에 대한 주도권을 다시 찾는 데 도움을 줍니다.}

\booksection{먼저 당신의 “화면 시간”을 보라——디지털이 사람을 맑게 함}

변경하기 전에, 먼저 현 상황을 파악. 스마트폰에는 많은 사람이 한 번도 열어본 적 없는 기능이 있는데, “화면 사용 시간”(iPhone) 또는 “디지털 건강”(안드로이드)이라 부르며, 매일 스마트폰에서 실제로 얼마나 시간을 쓰는지, 어떤 앱에 시간을 쓰는지 알려줌.

\begin{itemize}
  \item \textbf{iPhone}: 설정 > 화면 사용 시간 > 모든 활동 보기. 매일 잠금 해제 횟수, 스마트폰 집는 횟수, 각 앱 사용 시간을 볼 수 있음.
\end{itemize}

\textbf{처음 이 데이터를 보는 사람 중 80\%는 놀랍니다}. 제 친구 한 명이 본 후 자신은 매일 도우인(도우인)에 4시간 이상을 쓴다는 것을 발견, 일주일이면 28시간——이는 아르바이트의 시간과 맞먹음. 이러한 “디지털 충격”이 변화의 첫 걸음.

\booksection{앱에 “사용 한도” 설정——자신에게 속도 제한 장치 달기}

문제가 어디인지 알고 나면, 표적성 있게 제한을 설정 가능. 앱을 일괄 삭제할 필요 없이, 단지 “일일 한도”를 설정하면 됨:

\begin{enumerate}
  \item \textbf{iPhone}: 설정 > 화면 사용 시간 > 앱 한도 > 한도 추가. 제한할 앱 선택(예: 도우인), 매일 30분 설정. 시간 후 앱 아이콘이 회색으로 변하고, 열려면 추가 비밀번호 입력 필요.
  \item \textbf{안드로이드}: 설정 > 디지털 건강 > 앱 타이머. 앱 선택 후 일일 사용 상한 설정, 시간 후 앱 아이콘이 회색으로 변함.
  \item \textbf{소기술}: 한도 비밀번호를 자주 사용하지 않고, 잘 기억하기 어려운 것으로 설정. 너무 쉽게 입력하면, 매번 "수명 연장"하고 싶어져, 한도가 사실상 무용지물이 됨.
\end{enumerate}

\booksection{“방해 금지 모드”의 올바른 사용법——알림을 끄는 것이 아니라 불안을 끄는 것}

많은 사람들의 스마트폰은 하루에 수십 번 알림이 울리고, 매번 울릴 때마다 주의가 한 번씩 중단됨. “알림을 한 번 본다”고 생각하지만, 실제로는 매번 중단 후 주의력을 다시 집중하는 데 15-20분 소요. 방해 금지 모드는 “연락 끊기”가 아니라, 주의력을 다시 통제할 수 있도록 도움.

\begin{itemize}
  \item \textbf{정시 켜기}: 방해 금지 모드를 밤 10시\textasciitilde{}아침 7시에 자동 켜지도록 설정. 이렇게 하면 잠잘 때 스마트폰이 스팸 문자에 깨어나지 않음.
  \item \textbf{중요 연락처는 통과 허용}: 방해 금지 모드 설정에서, "전화 허용"을 "즐겨찾기 연락처" 또는 "반복 전화"(3분 내 두 번째 전화는 정상 울림)로 설정. 이렇게 하면 가족의 긴급 전화를 놓치지 않으면서, 영업 전화는 방해받지 않음.
  \item \textbf{작업 집중 모드}: iPhone의 "집중 모드"는 사용자 정의 시나리오를 설정 가능——예를 들어 "작업" 모드에서는 작업 관련 앱 알림만 허용, "독서" 모드에서는 모든 알림이 무음.
\end{itemize}

\booksection{눈 보호는 앱에 의존하지 말고 습관에 의존——20-20-20 법칙}

오랜 시간 스마트폰 화면을 보면, 눈이 건조하고 뻐근해지는 것은 필연. 눈 보호 모드(블루라이트 차단)는 어느 정도 도움이 되지만, 진짜 효과적인 것은 눈 사용 습관의 변화:

\begin{tcbinfo}{20-20-20 법칙}
매\textbf{20분}마다 시선을\textbf{20피트(약 6미터)}밖으로 옮기고,\textbf{20초}유지. 이 간단한 동작은 섬모체 피로를 효과적으로 완화.

\end{tcbinfo}

\begin{itemize}
  \item \textbf{눈 보호 모드 켜기}: 설정 > 디스플레이 > 눈 보호 모드(또는 나이트 시프트), 켠 후 화면 색온도가 따뜻해져 블루라이트 감소. 일몰에서 일출까지 자동 켜기로 설정 권장.
  \item \textbf{화면 밝기 조정}: 환경이 어두우면 화면도 어두워야 함. 많은 사람들이 어둠 속에서 스마트폰을 볼 때 밝기를 너무 높게 설정하는데, 이것이 눈을 가장 심하게 손상시키는 행동. "자동 밝기"를 켜고 정상 작동 확인.
\end{itemize}

\booksection{“스마트폰 중독족”의 목 자구 가이드}

스마트폰을 보면서 고개를 숙이면, 목이 받는 압력은 정상 기립의 5-6배. 장기간 고개를 숙이면 목 곡도가 곧게 변하고, 골질 증생, 어지러움, 손 저림이 발생. 다음은 간단한 개선 방법:

\begin{itemize}
  \item \textbf{스마트폰을 높이 들어올리기}: 스마트폰을 볼 때 가능한 한 눈과 같은 높이로 높이 들어올리기, 다리 위에 놓인 스마트폰을 고개를 숙여 보지 말기. 팔은 피곤하지만, 목이 감사할 것.
  \item \textbf{30분마다 목 움직이기}: 천천히 왼쪽으로 머리를 끝까지 돌리고, 3초 멈추고, 정면; 오른쪽으로 돌리고, 3초 멈추기. 5회 반복. 그리고 "미자체조"——턱으로 공중에 "미"자 쓰기.
  \item \textbf{벽에 기대어 서기 3분}: 뒤통수, 견갑골, 둔부, 발뒤꿈치 4점을 벽에 붙이고, 매일 3분 서기, 가슴 숙이고 고개 드리는 자세 교정에 효과적.
\end{itemize}

\booksection{수면과 스마트폰——침실에서 가장 고쳐야 할 습관}

스마트폰 화면의 블루라이트는 멜라토닌 분비를 억제하여, “볼수록 더 정신이 깨어나게” 함. 많은 사람들이 “보면서 졸린다”고 느끼지만, 실제로는 블루라이트에 의해 뇌가 속아, 스마트폰을 내려놓으면 오히려 잠들기 어려워짐. 다음은 개선 제안:

\begin{itemize}
  \item \textbf{취침 1시간 전 스마트폰 보지 않기}: 가장 효과적이지만 가장 지키기 어려운 규칙. 지키지 못하면, 적어도 "나이트 시프트/눈 보호 모드"를 켜 블루라이트를 최저로 낮춤.
  \item \textbf{스마트폰을 손 닿지 않는 곳에 두기}: 충전기를 침실 반대편에 두기, 침대 협탁이 아님. 이렇게 하면 참다못해 집어 보지도 않을 뿐더러, 알람이 울려도 일어나야 끄게——일석이조.
  \item \textbf{"수면 모드" 켜기}: iPhone의 "수면" 집중 모드는 수면 시간에 자동으로 화면 잠금, 잠금 화면 간소화; 안드로이드 각 브랜드에도 유사 기능.
  \item \textbf{알람 대신 알람시계 사용}: 1만 원대 실체 알람시계 구매, 스마트폰을 알람으로 사용하지 않기, "시간을 보려고 그리고 참다못해 보는" 함정 감소.
\end{itemize}

\booksection{정보 불안증——“선택적 실명” 배우기}

우리가 매일 받는 정보량은, 옛사람이 평생 받은 양. 그러나 인간의 뇌는 이렇게 많은 정보를 처리하도록 설계되지 않았으며, 정보 과부하가 불안, 주의력 산만, 결정 피로를 유발. 해결책은 “더 많은 정보 얻기”가 아니라 “능동적 차단”:

\begin{itemize}
  \item \textbf{팔로우 취소}: 팔로우하는 모든 공식 계정, UP주, 블로거 목록을 열고, 자신에게 질문: "지난 한 달 동안 실제로 가치를 얻은 적이 있나?" 없다면 취소. 정보 채널은 적지만 귀한 것이지 많은 것이 귀한 것이 아님.
  \item \textbf{불필요한 앱 푸시 끄기}: 위챗, 전화, 문자, 일정표 외의 다른 앱 알림을 모두 끄기. 뉴스는 푸시가 필요 없음——자발적으로 볼 때 자연히 보임.
  \item \textbf{"정보 단식일" 설정}: 매주 하루(예: 일요일)를 선택, 종일 스마트폰 뉴스 보지 않고, 숏폼 영상 보지 않기. 세상이 무너지지 않으며, 월요일 정신 상태가 눈에 띄게 좋아짐.
\end{itemize}

\begin{tcbpractice}{실습: 일주일 디지털 건강 챌린지}
\begin{enumerate}
  \item 오늘: 화면 사용 시간/디지털 건강 열고, 지난주 사용 데이터 확인.
  \item 내일: 가장 오래 사용한 엔터테인먼트 앱에 일일 한도 설정(30-60분 권장).
  \item 모레: 방해 금지 모드 정시 켜기(22:00-07:00), 즐겨찾기 연락처 전화 허용.
  \item 나흘째: 취침 1시간 전 스마트폰 보지 않기, 충전기를 침실 반대편으로.
  \item 다섯째 날: 더 이상 보지 않는 공식 계정 또는 블로거 5개 팔로우 취소.
  \item 여섯째 날: 하루 종일 "정보 단식" 시도——뉴스 안 보고 숏폼 영상 안 보기.
  \item 일곱째 날: 이번 주 화면 사용 시간 보고서 확인, 지난주와 비교.
\end{enumerate}

\end{tcbpractice}

\begin{tcbsummary}{본 챕터 요약}
\begin{itemize}
  \item 먼저 화면 사용 시간 데이터를 보고, 현 상황을 인식하는 것이 변화의 첫 걸음
  \item 엔터테인먼트 앱에 한도 설정, 방해 금지 모드 정시 설정, 주의력 다시 장악
  \item 20-20-20 법칙은 가장 간단하고 효과적인 눈 보호 습관
  \item 스마트폰 보면서 고개 숙이면 목에 5-6배 하중, 스마트폰 들어올리기+정기적으로 목 움직이기
  \item 취침 1시간 전 스마트폰 보지 않기, 충전기를 손 닿지 않는 곳에 두기
  \item 정보 불안의 해독제는 "능동적 차단", "더 얻기"가 아님
\end{itemize}

\end{tcbsummary}
